\chapter{Metodología experimental}
\label{ch:evaluacion}
% Fuentes canónicas: FINAL_HOLDOUT_EVALUATION_CONTRACT_V2.md;
% FINAL_HOLDOUT_V2_RESULTS.md; V2_B_EVALUATOR.md. No ejecutar evaluación.
\section{Diseño independiente y selección del holdout}
\pendiente{Retomar la separación del capítulo~\ref{ch:datos}: 48 BOE,
seis estratos, exclusión de exposición y selección previa. Precisar que se
evalúa la salida documental del sistema congelado, sin medir agrupación Gold
ni usabilidad. No generalizar al BOE completo.}
\pendiente{TABLA — Composición del holdout: 48 BOE en seis estratos
A/B por año, ocho por estrato. Explicar la separación respecto del corpus
de desarrollo y remitir a la composición del capítulo 3 para evitar duplicarla.}
\section{Referencia humana y anotación ciega}
% Contrato V2 §§3--9; resultados §§2,10; interfaz evaluation/final_holdout_v2/.
\pendiente{Definir ground truth como referencia anotada por una persona y
anotación ciega como trabajo sin consultar predicciones del sistema.
Describir campos primarios V2, interfaz auxiliar Streamlit y QA efectivamente
realizado. No inferir segunda revisión independiente ni acuerdo interanotador
por disponer de un exportador de QA o metadatos de revisión.}
\pendiente{CAPTURA C04, opcional — Interfaz de anotación ciega.
Solo si aclara el procedimiento: vista pequeña, legible y sin confundirla con
la aplicación pública. Generar después, con caption y referencia textual.}
\section{Versiones congeladas y orden del experimento}
% Resultados §§1,11; TFM_CLOSEOUT.md, checkpoint septiembre.
\pendiente{Explicar freeze como fijación de versiones antes de conocer
predicciones. Registrar sistema, truth V2 y evaluador definitivo V2-B,
diferenciando el evaluador histórico conservado. Separar selección del
holdout, cierre productivo y congelación preejecución.}
\pendiente{FIGURA F07 — Metodología experimental
\par Etiqueta prevista: \path{fig:metodologia-experimental}. Dos ramas: corpus de desarrollo
hacia sistema congelado; holdout independiente hacia anotación humana ciega
y truth V2. Mostrar freeze de truth y del evaluador definitivo V2-B antes
de la ejecución primaria; después, predicciones congeladas, evaluación,
resultados primarios y análisis post-hoc separado. Las flechas indican
dependencias, sin inventar fechas ni una independencia interanotador.
Usar exclusivamente la cronología de los resultados canónicos. Referenciar
desde esta sección; no repetir el diagrama en reproducibilidad.}
\section{Emparejamiento y reglas de puntuación V2-B}
% Contrato V2 §§11--12; evaluation/final_holdout_v2/scoring_rules.json.
\pendiente{Definir matching como correspondencia entre elementos de la
predicción y de la referencia. Explicar progresivamente nombres de activos,
con un ejemplo de relación nominal reconocible y coincidencia normalizada
exacta exigida por el evaluador; usar Hipódromo solo como ilustración del
diagnóstico ya documentado, sin crear una nueva correspondencia.
Continuar con el
conjunto completo de activos del evento, evidencia de actuaciones y
localizaciones dentro de eventos emparejados; correspondencias uno a uno.
Describir atribución actuación a activo y denominadores sin modificarlos.}
\section{Temporalidad, evidencia y salvaguarda P0}
\pendiente{Distinguir actuaciones actuales de antecedentes; estos últimos
no integran el denominador positivo de detección actual. Explicar evidencia
aceptada, fragmentos y adjudicación de P0. Antecedente sin extracción
emparejada no es un falso negativo de P0; no existe accuracy temporal
categórica. Mantener separado el ensayo de desarrollo.}
\section{Métricas y ejemplo didáctico}
\pendiente{Definir primero TP (correctamente detectado), FP (predicho sin
correspondencia) y FN (esperado sin detección); después precisión, recall,
F1 y accuracy. Ejemplo exclusivamente didáctico: diez elementos esperados,
ocho predichos y seis correctos; precisión 6/8, recall 6/10. Explicar que F1
es su media armónica, accuracy corresponde a clasificación y no existe una
accuracy global. Tratar denominadores nulos y atributos condicionados.
TODO-CITA: fuente metodológica de evaluación de extracción.}
\section{Reproducibilidad y registro de ejecución}
% Resultados §§1,7,11; PRIMARY_EXECUTION_V2.md.
\pendiente{Definir manifest, identidad de versión y hash antes de presentar
las identidades; remitir su detalle al anexo. Describir preservación de
predicciones, evaluador y registro de ejecución, incluyendo el error terminal.
No prometer identidad de nuevas respuestas generativas. Remitir a F07 para
las barreras de exposición, sin crear una segunda figura del mismo proceso.}
